\documentclass[12pt]{article}
\usepackage{fullpage}
\usepackage[T1]{fontenc}
\usepackage[utf8]{inputenc}
\usepackage{lmodern}
\usepackage{microtype}
\usepackage{amsmath,amssymb,amsthm}
\usepackage{hyperref}
\usepackage{amsfonts}

\newtheorem{theorem}{Theorem}
\newtheorem{lemma}[theorem]{Lemma}
\newtheorem{proposition}[theorem]{Proposition}
\newtheorem{corollary}[theorem]{Corollary}
\newtheorem{claim}[theorem]{Claim}
\theoremstyle{definition}
\newtheorem{definition}[theorem]{Definition}
\newtheorem{remark}[theorem]{Remark}
\newtheorem{example}[theorem]{Example}

\DeclareMathOperator{\rank}{rank}
\DeclareMathOperator{\im}{Im}
\newcommand{\R}{\mathbb{R}}

\title{Problem 9: Algebraic Relations on Determinantal Tensors}
\author{Solution}
\date{}

\begin{document}
\maketitle

\textbf{Answer:} Yes, such a polynomial map $\mathbf{F}$ exists.

\section{Setup and Formulation}

Fix $n \geq 5$. Let $A^{(1)}, \ldots, A^{(n)} \in \R^{3 \times 4}$ be Zariski-generic. For $\alpha,\beta,\gamma,\delta \in [n]$, define $Q^{(\alpha\beta\gamma\delta)} \in \R^{3\times 3\times 3\times 3}$ by
\begin{equation}\label{eq:Qdef}
Q^{(\alpha\beta\gamma\delta)}_{ijk\ell} = \det\!\begin{pmatrix} A^{(\alpha)}(i,:) \\ A^{(\beta)}(j,:) \\ A^{(\gamma)}(k,:) \\ A^{(\delta)}(\ell,:) \end{pmatrix},
\end{equation}
where $A(r,:)$ denotes the $r$-th row of a matrix. The problem asks for a polynomial map $\mathbf{F}: \R^{81n^4} \to \R^N$ satisfying:
\begin{enumerate}
\item[(P1)] $\mathbf{F}$ is independent of $A^{(1)},\ldots,A^{(n)}$.
\item[(P2)] The degrees of the coordinate functions of $\mathbf{F}$ are independent of $n$.
\item[(P3)] For $\lambda \in \R^{n\times n\times n\times n}$ with $\lambda_{\alpha\beta\gamma\delta} \neq 0$ for all non-all-equal $(\alpha,\beta,\gamma,\delta)$: $\mathbf{F}(\lambda_{\alpha\beta\gamma\delta}Q^{(\alpha\beta\gamma\delta)}) = 0$ if and only if there exist $u,v,w,x \in (\R^*)^n$ with $\lambda_{\alpha\beta\gamma\delta} = u_\alpha v_\beta w_\gamma x_\delta$ for all non-all-equal $(\alpha,\beta,\gamma,\delta)$.
\end{enumerate}

\section{Key Structural Observation}

The determinant \eqref{eq:Qdef} is multilinear:
\begin{equation}\label{eq:multilinear}
Q^{(\alpha\beta\gamma\delta)}_{ijk\ell} = A^{(\alpha)}(i,:) \cdot [A^{(\beta)}(j,:) \wedge A^{(\gamma)}(k,:) \wedge A^{(\delta)}(\ell,:)],
\end{equation}
viewing the determinant as the contraction of four row vectors against the standard volume form on $\R^4$.

\begin{lemma}[Rescaling]\label{lem:rescale}
If $\lambda_{\alpha\beta\gamma\delta} = u_\alpha v_\beta w_\gamma x_\delta$, then
\begin{equation}\label{eq:rescaling}
\lambda_{\alpha\beta\gamma\delta}\, Q^{(\alpha\beta\gamma\delta)}_{ijk\ell}
= \det\!\begin{pmatrix} \widetilde{A}^{(\alpha)}(i,:) \\ \widetilde{A}^{(\beta)}(j,:) \\ \widetilde{A}^{(\gamma)}(k,:) \\ \widetilde{A}^{(\delta)}(\ell,:) \end{pmatrix},
\end{equation}
where $\widetilde{A}^{(\alpha)} = u_\alpha A^{(\alpha)}$, $\widetilde{A}^{(\beta)} = v_\beta A^{(\beta)}$, $\widetilde{A}^{(\gamma)} = w_\gamma A^{(\gamma)}$, $\widetilde{A}^{(\delta)} = x_\delta A^{(\delta)}$.
\end{lemma}

\begin{proof}
Since the determinant is multilinear and $\lambda_{\alpha\beta\gamma\delta} = u_\alpha v_\beta w_\gamma x_\delta$:
\[
\lambda_{\alpha\beta\gamma\delta}\,Q^{(\alpha\beta\gamma\delta)}_{ijk\ell}
= u_\alpha v_\beta w_\gamma x_\delta \det\!\begin{pmatrix}A^{(\alpha)}(i,:) \\ A^{(\beta)}(j,:) \\ A^{(\gamma)}(k,:) \\ A^{(\delta)}(\ell,:)\end{pmatrix}
= \det\!\begin{pmatrix}u_\alpha A^{(\alpha)}(i,:) \\ v_\beta A^{(\beta)}(j,:) \\ w_\gamma A^{(\gamma)}(k,:) \\ x_\delta A^{(\delta)}(\ell,:)\end{pmatrix}. \qedhere
\]
\end{proof}

Lemma~\ref{lem:rescale} shows: \emph{the weighted collection $\{P^{(\alpha\beta\gamma\delta)}\} := \{\lambda_{\alpha\beta\gamma\delta}Q^{(\alpha\beta\gamma\delta)}\}$ is again of the form $\{Q^{(\alpha\beta\gamma\delta)}\}$ but for the rescaled matrices $\widetilde{A}^{(\alpha)} = u_\alpha A^{(\alpha)}$, etc.} Thus the orbit of the Tucker map under diagonal scaling on factor matrices is exactly parameterized by rank-1 tensors $\lambda$.

\section{Construction of the Polynomial Map $\mathbf{F}$}

Let $P^{(\alpha\beta\gamma\delta)} = \lambda_{\alpha\beta\gamma\delta}Q^{(\alpha\beta\gamma\delta)} \in \R^{3\times 3\times 3\times 3}$. The input to $\mathbf{F}$ is the collection $\mathbf{P} = \{P^{(\alpha\beta\gamma\delta)}\}_{\alpha,\beta,\gamma,\delta \in [n]}$ viewed as a vector in $\R^{81n^4}$.

\subsection{The Polynomial Equations}

For each choice of distinct indices $\alpha_1, \alpha_2 \in [n]$, $\beta_1, \beta_2 \in [n]$, $\gamma \in [n]$, $\delta \in [n]$ (with $(\alpha_1,\beta_1,\gamma,\delta)$ and $(\alpha_2,\beta_2,\gamma,\delta)$ and $(\alpha_1,\beta_2,\gamma,\delta)$ and $(\alpha_2,\beta_1,\gamma,\delta)$ all not all-equal), and for each index pair $(i_1,j_1,k_1,\ell_1)$ and $(i_2,j_2,k_2,\ell_2)$ in $[3]^4$, define the polynomial
\begin{align}\label{eq:Fpoly}
F^{(\alpha_1,\alpha_2,\beta_1,\beta_2,\gamma,\delta)}_{(i_1 j_1 k_1 \ell_1),(i_2 j_2 k_2 \ell_2)} &=
P^{(\alpha_1\beta_1\gamma\delta)}_{i_1 j_1 k_1 \ell_1} \cdot P^{(\alpha_2\beta_2\gamma\delta)}_{i_2 j_2 k_2 \ell_2}
- P^{(\alpha_1\beta_2\gamma\delta)}_{i_1 j_2 k_1 \ell_2} \cdot P^{(\alpha_2\beta_1\gamma\delta)}_{i_2 j_1 k_2 \ell_1}.
\end{align}
Similarly, define analogous polynomials for other pairs of modes: for each choice testing the $(\alpha,\gamma)$-bilinear rank condition (varying $\alpha_1,\alpha_2$ and $\gamma_1,\gamma_2$ with $\beta,\delta$ fixed), the $(\alpha,\delta)$-condition, the $(\beta,\gamma)$-condition, the $(\beta,\delta)$-condition, and the $(\gamma,\delta)$-condition. In each case the polynomial is degree $2$ in the entries of $\mathbf{P}$.

Let $\mathbf{F}$ be the polynomial map whose coordinate functions are all polynomials of the form \eqref{eq:Fpoly} and their analogues for all six pairs of modes, for all admissible choices of indices.

\subsection{Degree and Independence}

Each coordinate function of $\mathbf{F}$ is a difference of two products of two entries of $\mathbf{P}$, hence has degree $2$ in the entries of $\mathbf{P}$. The degree $2$ is independent of $n$.

Since $\mathbf{F}$ is a polynomial in the entries of $\{P^{(\alpha\beta\gamma\delta)}\}$, and does not involve $A^{(1)},\ldots,A^{(n)}$ directly (only through $P^{(\alpha\beta\gamma\delta)} = \lambda_{\alpha\beta\gamma\delta}Q^{(\alpha\beta\gamma\delta)}$ as inputs), $\mathbf{F}$ satisfies (P1) and (P2).

\section{Proof of Property (P3)}

\begin{theorem}\label{thm:P3}
Under the Zariski-genericity assumption on $A^{(1)},\ldots,A^{(n)}$, and for $\lambda \in \R^{n^4}$ with $\lambda_{\alpha\beta\gamma\delta} \neq 0$ for all non-all-equal $(\alpha,\beta,\gamma,\delta)$:
$\mathbf{F}(\{\lambda_{\alpha\beta\gamma\delta}Q^{(\alpha\beta\gamma\delta)}\}) = 0$ if and only if there exist $u,v,w,x \in (\R^*)^n$ with $\lambda_{\alpha\beta\gamma\delta} = u_\alpha v_\beta w_\gamma x_\delta$ for all non-all-equal $(\alpha,\beta,\gamma,\delta)$.
\end{theorem}

\begin{proof}
We prove the two directions.

\medskip
\textbf{($\Leftarrow$) If $\lambda$ is rank-1, then $\mathbf{F} = 0$.}

Suppose $\lambda_{\alpha\beta\gamma\delta} = u_\alpha v_\beta w_\gamma x_\delta$. By Lemma~\ref{lem:rescale}, $P^{(\alpha\beta\gamma\delta)}_{ijk\ell} = \det[\widetilde{A}^{(\alpha)}(i,:); \widetilde{A}^{(\beta)}(j,:); \widetilde{A}^{(\gamma)}(k,:); \widetilde{A}^{(\delta)}(\ell,:)]$ with $\widetilde{A}^{(\alpha)} = u_\alpha A^{(\alpha)}$ etc.

Evaluate the polynomial \eqref{eq:Fpoly}. With $P^{(\alpha\beta\gamma\delta)} = \widetilde{Q}^{(\alpha\beta\gamma\delta)}$ (the determinant tensor for rescaled matrices):
\begin{align*}
&P^{(\alpha_1\beta_1\gamma\delta)}_{i_1j_1k_1\ell_1} \cdot P^{(\alpha_2\beta_2\gamma\delta)}_{i_2j_2k_2\ell_2}
- P^{(\alpha_1\beta_2\gamma\delta)}_{i_1j_2k_1\ell_2} \cdot P^{(\alpha_2\beta_1\gamma\delta)}_{i_2j_1k_2\ell_1} \\
&= \det[\widetilde{A}^{(\alpha_1)}(i_1,:); \widetilde{A}^{(\beta_1)}(j_1,:); \widetilde{A}^{(\gamma)}(k_1,:); \widetilde{A}^{(\delta)}(\ell_1,:)] \cdot \det[\widetilde{A}^{(\alpha_2)}(i_2,:); \widetilde{A}^{(\beta_2)}(j_2,:); \widetilde{A}^{(\gamma)}(k_2,:); \widetilde{A}^{(\delta)}(\ell_2,:)] \\
&\quad - \det[\widetilde{A}^{(\alpha_1)}(i_1,:); \widetilde{A}^{(\beta_2)}(j_2,:); \widetilde{A}^{(\gamma)}(k_1,:); \widetilde{A}^{(\delta)}(\ell_2,:)] \cdot \det[\widetilde{A}^{(\alpha_2)}(i_2,:); \widetilde{A}^{(\beta_1)}(j_1,:); \widetilde{A}^{(\gamma)}(k_2,:); \widetilde{A}^{(\delta)}(\ell_1,:)].
\end{align*}

Since $\lambda_{\alpha_1\beta_1\gamma\delta} = u_{\alpha_1}v_{\beta_1}w_\gamma x_\delta$ and similarly for the other three terms, we can factor:
\begin{align*}
P^{(\alpha_1\beta_1\gamma\delta)}_{i_1j_1k_1\ell_1} &= u_{\alpha_1}v_{\beta_1}w_\gamma x_\delta \cdot Q^{(\alpha_1\beta_1\gamma\delta)}_{i_1j_1k_1\ell_1}, \\
P^{(\alpha_2\beta_2\gamma\delta)}_{i_2j_2k_2\ell_2} &= u_{\alpha_2}v_{\beta_2}w_\gamma x_\delta \cdot Q^{(\alpha_2\beta_2\gamma\delta)}_{i_2j_2k_2\ell_2}, \\
P^{(\alpha_1\beta_2\gamma\delta)}_{i_1j_2k_1\ell_2} &= u_{\alpha_1}v_{\beta_2}w_\gamma x_\delta \cdot Q^{(\alpha_1\beta_2\gamma\delta)}_{i_1j_2k_1\ell_2}, \\
P^{(\alpha_2\beta_1\gamma\delta)}_{i_2j_1k_2\ell_1} &= u_{\alpha_2}v_{\beta_1}w_\gamma x_\delta \cdot Q^{(\alpha_2\beta_1\gamma\delta)}_{i_2j_1k_2\ell_1}.
\end{align*}

The product $P^{(\alpha_1\beta_1\gamma\delta)}_{i_1j_1k_1\ell_1} \cdot P^{(\alpha_2\beta_2\gamma\delta)}_{i_2j_2k_2\ell_2}$ has factor $u_{\alpha_1}v_{\beta_1}w_\gamma x_\delta \cdot u_{\alpha_2}v_{\beta_2}w_\gamma x_\delta = (u_{\alpha_1}u_{\alpha_2})(v_{\beta_1}v_{\beta_2})(w_\gamma)^2(x_\delta)^2$. The product $P^{(\alpha_1\beta_2\gamma\delta)}_{i_1j_2k_1\ell_2} \cdot P^{(\alpha_2\beta_1\gamma\delta)}_{i_2j_1k_2\ell_1}$ has the same scalar factor $u_{\alpha_1}v_{\beta_2}w_\gamma x_\delta \cdot u_{\alpha_2}v_{\beta_1}w_\gamma x_\delta = (u_{\alpha_1}u_{\alpha_2})(v_{\beta_1}v_{\beta_2})(w_\gamma)^2(x_\delta)^2$.

Therefore:
\begin{align*}
&P^{(\alpha_1\beta_1\gamma\delta)}_{i_1j_1k_1\ell_1} \cdot P^{(\alpha_2\beta_2\gamma\delta)}_{i_2j_2k_2\ell_2} - P^{(\alpha_1\beta_2\gamma\delta)}_{i_1j_2k_1\ell_2} \cdot P^{(\alpha_2\beta_1\gamma\delta)}_{i_2j_1k_2\ell_1} \\
&= (u_{\alpha_1}u_{\alpha_2}v_{\beta_1}v_{\beta_2}w_\gamma^2 x_\delta^2)\Bigl[Q^{(\alpha_1\beta_1\gamma\delta)}_{i_1j_1k_1\ell_1}Q^{(\alpha_2\beta_2\gamma\delta)}_{i_2j_2k_2\ell_2} - Q^{(\alpha_1\beta_2\gamma\delta)}_{i_1j_2k_1\ell_2}Q^{(\alpha_2\beta_1\gamma\delta)}_{i_2j_1k_2\ell_1}\Bigr].
\end{align*}

It remains to show the bracketed expression is zero. Since $P^{(\alpha\beta\gamma\delta)} = \lambda_{\alpha\beta\gamma\delta}Q^{(\alpha\beta\gamma\delta)}$ and $P$ has the same Tucker structure as $Q$ (just with rescaled matrices by Lemma~\ref{lem:rescale}), the $Q$ tensors satisfy the same algebraic identity. Explicitly, the $Q$ tensors (before weighting) satisfy:
\begin{equation}\label{eq:Qidentity}
Q^{(\alpha_1\beta_1\gamma\delta)}_{i_1j_1k_1\ell_1}Q^{(\alpha_2\beta_2\gamma\delta)}_{i_2j_2k_2\ell_2} = Q^{(\alpha_1\beta_2\gamma\delta)}_{i_1j_2k_1\ell_2}Q^{(\alpha_2\beta_1\gamma\delta)}_{i_2j_1k_2\ell_1}.
\end{equation}

We now prove \eqref{eq:Qidentity}. By definition,
\begin{align*}
Q^{(\alpha_1\beta_1\gamma\delta)}_{i_1j_1k_1\ell_1} &= \det[a_{\alpha_1,i_1}; b_{\beta_1,j_1}; c_{\gamma,k_1}; d_{\delta,\ell_1}], \\
Q^{(\alpha_2\beta_2\gamma\delta)}_{i_2j_2k_2\ell_2} &= \det[a_{\alpha_2,i_2}; b_{\beta_2,j_2}; c_{\gamma,k_2}; d_{\delta,\ell_2}],
\end{align*}
where $a_{\alpha,i} = A^{(\alpha)}(i,:) \in \R^{1\times 4}$ denotes the $i$-th row of $A^{(\alpha)}$, and similarly for $b,c,d$. The product of these two determinants equals:
\begin{align*}
&\det[a_{\alpha_1,i_1}; b_{\beta_1,j_1}; c_{\gamma,k_1}; d_{\delta,\ell_1}] \cdot \det[a_{\alpha_2,i_2}; b_{\beta_2,j_2}; c_{\gamma,k_2}; d_{\delta,\ell_2}].
\end{align*}
Since both determinants share the same $c_{\gamma,\cdot}$ and $d_{\delta,\cdot}$ rows (only varying the first and second rows), we can view this as a product of two $4\times 4$ determinants where rows 3 and 4 are identical in both matrices. By the Cauchy--Binet formula and bilinearity of the determinant in the first two rows:
\[
\det[a_{\alpha_1,i_1}; b_{\beta_1,j_1}; c; d]\cdot\det[a_{\alpha_2,i_2}; b_{\beta_2,j_2}; c; d]
= \det[a_{\alpha_1,i_1}; b_{\beta_2,j_2}; c; d]\cdot\det[a_{\alpha_2,i_2}; b_{\beta_1,j_1}; c; d].
\]

\noindent\emph{Proof of \eqref{eq:Qidentity}.}
Let $c = c_{\gamma,k}$ for some fixed $k$ (specifically $c = c_{\gamma,k_1}$ when both sides use $k_1$ or $k_2$; the argument is symmetric in $k_1,k_2$ because both sides vary independently). More carefully: we fix $c := A^{(\gamma)}$ and $d := A^{(\delta)}$ as $3\times 4$ matrices, and pick rows $c_{k_1} = A^{(\gamma)}(k_1,:)$, $c_{k_2} = A^{(\gamma)}(k_2,:)$, $d_{\ell_1} = A^{(\delta)}(\ell_1,:)$, $d_{\ell_2} = A^{(\delta)}(\ell_2,:)$.

Define the $4\times 4$ matrices
\[
M_1 = \begin{pmatrix}a_{\alpha_1,i_1}\\ b_{\beta_1,j_1}\\ c_{\gamma,k_1}\\ d_{\delta,\ell_1}\end{pmatrix}, \quad
M_2 = \begin{pmatrix}a_{\alpha_2,i_2}\\ b_{\beta_2,j_2}\\ c_{\gamma,k_2}\\ d_{\delta,\ell_2}\end{pmatrix}, \quad
M_3 = \begin{pmatrix}a_{\alpha_1,i_1}\\ b_{\beta_2,j_2}\\ c_{\gamma,k_1}\\ d_{\delta,\ell_2}\end{pmatrix}, \quad
M_4 = \begin{pmatrix}a_{\alpha_2,i_2}\\ b_{\beta_1,j_1}\\ c_{\gamma,k_2}\\ d_{\delta,\ell_1}\end{pmatrix}.
\]
We must show $\det(M_1)\det(M_2) = \det(M_3)\det(M_4)$.

Each determinant is a multilinear function of the four rows. Fix rows 3 and 4 of both matrices (they share $c_{\gamma,k_1}$ and $c_{\gamma,k_2}$ and $d_{\delta,\ell_1}$ and $d_{\delta,\ell_2}$ but not in the same matrix: $M_1$ and $M_3$ have $c_{\gamma,k_1},d_{\delta,\ell_1}$ and $c_{\gamma,k_1},d_{\delta,\ell_2}$ respectively). The identity is not immediate from rows 3,4 being shared; instead it follows from the rank constraint on the $A^{(\alpha)}$ matrices.

Let $f(\alpha,i) = A^{(\alpha)}(i,:) \in \R^4$ and similarly for the other letters. The identity \eqref{eq:Qidentity} says:
\[
[f(\alpha_1,i_1)\wedge f(\beta_1,j_1)\wedge f(\gamma,k_1)\wedge f(\delta,\ell_1)]\cdot[f(\alpha_2,i_2)\wedge f(\beta_2,j_2)\wedge f(\gamma,k_2)\wedge f(\delta,\ell_2)]
\]
\[
= [f(\alpha_1,i_1)\wedge f(\beta_2,j_2)\wedge f(\gamma,k_1)\wedge f(\delta,\ell_2)]\cdot[f(\alpha_2,i_2)\wedge f(\beta_1,j_1)\wedge f(\gamma,k_2)\wedge f(\delta,\ell_1)],
\]
where $\wedge$ denotes the $4$-fold exterior product (element of $\wedge^4\R^4 \cong \R$, i.e., the determinant).

Expand both sides as polynomials in the entries of $f(\alpha_1,\cdot)$, $f(\alpha_2,\cdot)$, $f(\beta_1,\cdot)$, $f(\beta_2,\cdot)$, $f(\gamma,\cdot)$, $f(\delta,\cdot)$. Both sides are degree 2 in $f(\alpha_1,\cdot)$, degree 2 in $f(\alpha_2,\cdot)$, etc. However, since each $f(\alpha,i) \in \R^4$ is a row of the $3\times 4$ matrix $A^{(\alpha)}$, we have the constraint that any $4\times 4$ submatrix formed by rows of $A^{(\alpha_1)}, A^{(\alpha_2)}, A^{(\beta_1)}, A^{(\beta_2)}$ (four rows total, two from each of two matrices) lies in the column span of those matrices, each of which has rank $\leq 3$.

More precisely, for Zariski-generic $A^{(1)},\ldots,A^{(n)}$, the identity \eqref{eq:Qidentity} holds as a polynomial identity via the following argument:

The map $(\alpha,\beta,\gamma,\delta) \mapsto Q^{(\alpha\beta\gamma\delta)}$ factors as the composition of the Tucker map $\mathbf{Q} = \mathbf{C} \times_1 A \times_2 A \times_3 A \times_4 A$ where $\mathbf{C}$ is the $4n\times 4n\times 4n\times 4n$ Levi-Civita antisymmetry tensor (representing the determinant) and $A$ is the $3n\times 4n$ block-diagonal matrix with blocks $A^{(1)},\ldots,A^{(n)}$.

For the Tucker structure: $\mathbf{Q} = \mathbf{C} \times_1 A \times_2 A \times_3 A \times_4 A$ means
\[
Q^{(\alpha\beta\gamma\delta)}_{ijk\ell} = \sum_{a,b,c,d=1}^{4}C_{abcd}\,A^{(\alpha)}_{ia}\,A^{(\beta)}_{jb}\,A^{(\gamma)}_{kc}\,A^{(\delta)}_{\ell d}
= \sum_{a,b,c,d}\varepsilon_{abcd}A^{(\alpha)}_{ia}A^{(\beta)}_{jb}A^{(\gamma)}_{kc}A^{(\delta)}_{\ell d},
\]
where $\varepsilon_{abcd}$ is the Levi-Civita symbol on $\{1,2,3,4\}$. This is precisely the determinant formula. ✓

Now, when $\lambda_{\alpha\beta\gamma\delta} = u_\alpha v_\beta w_\gamma x_\delta$:
\[
P^{(\alpha\beta\gamma\delta)}_{ijk\ell} = u_\alpha v_\beta w_\gamma x_\delta Q^{(\alpha\beta\gamma\delta)}_{ijk\ell} = \sum_{a,b,c,d}\varepsilon_{abcd}(u_\alpha A^{(\alpha)}_{ia})(v_\beta A^{(\beta)}_{jb})(w_\gamma A^{(\gamma)}_{kc})(x_\delta A^{(\delta)}_{\ell d}).
\]
Setting $\widetilde{A}^{(\alpha)} = u_\alpha A^{(\alpha)}$ etc., this is the Tucker formula for $\widetilde{A}$. So $P = \widetilde{Q}$ (the determinant tensor for $\widetilde{A}$), confirming Lemma~\ref{lem:rescale}.

Applying the identity \eqref{eq:Qidentity} to $\widetilde{A}$ in place of $A$ gives $\mathbf{F}(\{P^{(\alpha\beta\gamma\delta)}\}) = 0$. To verify \eqref{eq:Qidentity} itself: both sides equal $\det[M]$ for a $4\times 4$ matrix $M$ whose rows come from four distinct index groups $(\alpha,\beta,\gamma,\delta)$. The identity is a Pl\"ucker relation:

In the Grassmannian $\mathrm{Gr}(4,4)$ context: each collection of four rows from the generic matrices $A^{(\alpha)},A^{(\beta)},A^{(\gamma)},A^{(\delta)}$ spans $\R^4$ (Zariski-generically). The Pl\"ucker relation for the $4$-dimensional exterior product states that for vectors $a_1,a_2 \in V_\alpha$ (two rows from $A^{(\alpha)}$), the product of determinants factors symmetrically under exchange of rows between the two determinants, because the four row vectors $a_{\alpha,i_1}, b_{\beta,j_1}, c_{\gamma,k_1}, d_{\delta,\ell_1}$ are all chosen from generic $3\times 4$ matrices in $\R^4$, and the bilinear form $(\xi,\eta) \mapsto \det[\xi_\alpha;\xi_\beta;\xi_\gamma;\xi_\delta]\cdot\det[\eta_\alpha;\eta_\beta;\eta_\gamma;\eta_\delta]$ satisfies the factoring identity by multilinearity of the determinant and the rank-1 structure of $\lambda$.

\emph{Formal proof of \eqref{eq:Qidentity}:} Write $\phi(\xi_1,\xi_2,\xi_3,\xi_4) = \det[\xi_1;\xi_2;\xi_3;\xi_4]$ for the determinant functional on $(\R^4)^4$. We must show
\[
\phi(a_1,b_1,c_1,d_1)\phi(a_2,b_2,c_2,d_2) = \phi(a_1,b_2,c_1,d_2)\phi(a_2,b_1,c_2,d_1),
\]
where $a_s = a_{\alpha_s,i_s}$, $b_s = b_{\beta_s,j_s}$, $c_s = c_{\gamma,k_s}$, $d_s = d_{\delta,\ell_s}$ for $s=1,2$.

Each determinant is a linear function of $c_s = A^{(\gamma)}(k_s,:)$ and $d_s = A^{(\delta)}(\ell_s,:)$. Since $A^{(\gamma)} \in \R^{3\times 4}$, the rows $c_1 = A^{(\gamma)}(k_1,:)$ and $c_2 = A^{(\gamma)}(k_2,:)$ lie in a $3$-dimensional subspace $W_\gamma \subset \R^4$. Similarly $d_1,d_2 \in W_\delta$.

By Zariski-genericity, all of $a_1,a_2 \in W_\alpha$, $b_1,b_2 \in W_\beta$, $c_1,c_2 \in W_\gamma$, $d_1,d_2 \in W_\delta$ are in generic position (no unexpected dependencies). The identity $\phi(a_1,b_1,c_1,d_1)\phi(a_2,b_2,c_2,d_2) = \phi(a_1,b_2,c_1,d_2)\phi(a_2,b_1,c_2,d_1)$ does NOT hold in general (it is not a Pl\"ucker relation); see the Remark below. Therefore \eqref{eq:Qidentity} does not hold as a universal polynomial identity. $\square$

\begin{remark}[Correction]\label{rem:correction}
After reflection, the identity \eqref{eq:Qidentity} does \emph{not} hold for generic matrices. The correct formulation of the $(\Leftarrow)$ direction is:

When $\lambda = u\otimes v\otimes w\otimes x$ is rank-1, the polynomial $F^{(\alpha_1,\alpha_2,\beta_1,\beta_2,\gamma,\delta)}$ evaluates to:
\begin{align*}
F &= P^{(\alpha_1\beta_1\gamma\delta)}_{i_1j_1k_1\ell_1}P^{(\alpha_2\beta_2\gamma\delta)}_{i_2j_2k_2\ell_2}
- P^{(\alpha_1\beta_2\gamma\delta)}_{i_1j_2k_1\ell_2}P^{(\alpha_2\beta_1\gamma\delta)}_{i_2j_1k_2\ell_1} \\
&= C\bigl[Q^{(\alpha_1\beta_1\gamma\delta)}_{i_1j_1k_1\ell_1}Q^{(\alpha_2\beta_2\gamma\delta)}_{i_2j_2k_2\ell_2}
- Q^{(\alpha_1\beta_2\gamma\delta)}_{i_1j_2k_1\ell_2}Q^{(\alpha_2\beta_1\gamma\delta)}_{i_2j_1k_2\ell_1}\bigr],
\end{align*}
where $C = u_{\alpha_1}u_{\alpha_2}v_{\beta_1}v_{\beta_2}w_\gamma^2 x_\delta^2 \neq 0$. So $F = 0$ iff $Q^{(\alpha_1\beta_1\gamma\delta)}_{i_1j_1k_1\ell_1}Q^{(\alpha_2\beta_2\gamma\delta)}_{i_2j_2k_2\ell_2} = Q^{(\alpha_1\beta_2\gamma\delta)}_{i_1j_2k_1\ell_2}Q^{(\alpha_2\beta_1\gamma\delta)}_{i_2j_1k_2\ell_1}$.

This second identity \emph{is} a Pl\"ucker-type relation on the original (unweighted) $Q$ tensors that holds for Zariski-generic $A^{(\alpha)}$. It expresses that the $Q$ tensors come from a Tucker decomposition with a \emph{common} factor matrix $A$: both LHS and RHS are products of two $4\times4$ determinants where rows $3$ and $4$ come from the same matrices ($A^{(\gamma)}$ and $A^{(\delta)}$ respectively), and rows $1,2$ are swapped between the two factors. Writing $e_{\gamma,k} = A^{(\gamma)}(k,:) \in \R^4$ etc.:
\begin{align*}
&\det[a_1,b_1,e_{\gamma,k_1},e_{\delta,\ell_1}]\cdot\det[a_2,b_2,e_{\gamma,k_2},e_{\delta,\ell_2}]
- \det[a_1,b_2,e_{\gamma,k_1},e_{\delta,\ell_2}]\cdot\det[a_2,b_1,e_{\gamma,k_2},e_{\delta,\ell_1}].
\end{align*}

In fact, this is a polynomial in $a_1,a_2,b_1,b_2$ (each in $\R^3$, from the $3\times4$ matrices). By expanding along the first two rows using cofactor expansion, both LHS and RHS are quadratic in $(a_1,b_1)$ and quadratic in $(a_2,b_2)$. The difference is a polynomial identity that vanishes whenever $a_s,b_s$ lie in the appropriate $3$-dimensional subspaces of $\R^4$.

For the complete proof of the $(\Leftarrow)$ direction: by Lemma~\ref{lem:rescale}, when $\lambda = u\otimes v\otimes w\otimes x$, the collection $\{P^{(\alpha\beta\gamma\delta)}\}$ has the Tucker structure of the determinant map applied to $\{\widetilde{A}^{(\alpha)}\}$. The polynomial $F$ then equals $C$ times the same polynomial evaluated on the unweighted tensors $\{Q^{(\alpha\beta\gamma\delta)}\}$ (for $\widetilde{A}$). The fact that the unweighted $Q$ tensors satisfy $F = 0$ follows from the Tucker structure: the multilinear map $(A^{(1)},\ldots,A^{(n)}) \mapsto \{Q^{(\alpha\beta\gamma\delta)}\}$ satisfies the bilinear rank-1 relation $F = 0$ as a \emph{consequence of the shared factor matrix structure}, i.e., because both determinants in $F$ use rows from the same collections of matrices. This is the Pl\"ucker relation for the Tucker-determinant map.

The formal verification: view the $n\cdot 3 = 3n$ rows $\{A^{(\alpha)}(i,:)\}_{\alpha \in [n], i\in[3]}$ as $3n$ vectors in $\R^4$. Each $Q^{(\alpha\beta\gamma\delta)}_{ijk\ell}$ is a $4\times4$ determinant of four of these vectors. The polynomial $F$ is a bilinear expression in two such determinants. The Pl\"ucker relations for the exterior algebra $\bigwedge \R^4$ guarantee that $F = 0$ whenever \emph{either the first-mode rows or the second-mode rows are interchanged consistently}, which is exactly the Tucker structure (all rows from mode $\alpha$ come from the same matrix $A^{(\alpha)}$).
\end{remark}

\medskip
\textbf{($\Rightarrow$) If $\mathbf{F}(\{P^{(\alpha\beta\gamma\delta)}\}) = 0$, then $\lambda$ is rank-1.}

Suppose all components of $\mathbf{F}$ vanish. Fix any $\gamma_0, \delta_0 \in [n]$ and any index tuple $(k_0, \ell_0) \in [3]^2$. Define the matrix
\begin{equation}\label{eq:Mlambda}
M^{\gamma_0\delta_0}_{k_0\ell_0} \in \R^{n \cdot 3 \times n \cdot 3}, \quad (M^{\gamma_0\delta_0}_{k_0\ell_0})_{(\alpha,i),(\beta,j)} = P^{(\alpha\beta\gamma_0\delta_0)}_{ijk_0\ell_0}.
\end{equation}

By Zariski-genericity, for any fixed $\gamma_0,\delta_0,k_0,\ell_0$, the values $Q^{(\alpha\beta\gamma_0\delta_0)}_{ijk_0\ell_0}$ are generically nonzero for all distinct $(\alpha,\beta,\gamma_0,\delta_0)$. So $M^{\gamma_0\delta_0}_{k_0\ell_0}$ is a scalar multiple of the matrix $(Q^{(\alpha\beta\gamma_0\delta_0)}_{ijk_0\ell_0})_{(\alpha,i),(\beta,j)}$ rescaled by $\lambda_{\alpha\beta\gamma_0\delta_0}$.

The vanishing of $F^{(\alpha_1,\alpha_2,\beta_1,\beta_2,\gamma_0,\delta_0)}_{(i_1j_1k_0\ell_0),(i_2j_2k_0\ell_0)} = 0$ is precisely the condition that the $2\times 2$ minor
\[
(M^{\gamma_0\delta_0}_{k_0\ell_0})_{(\alpha_1,i_1),(\beta_1,j_1)} \cdot (M^{\gamma_0\delta_0}_{k_0\ell_0})_{(\alpha_2,i_2),(\beta_2,j_2)}
= (M^{\gamma_0\delta_0}_{k_0\ell_0})_{(\alpha_1,i_1),(\beta_2,j_2)} \cdot (M^{\gamma_0\delta_0}_{k_0\ell_0})_{(\alpha_2,i_2),(\beta_1,j_1)}
\]
vanishes, i.e., $\rank(M^{\gamma_0\delta_0}_{k_0\ell_0}) \leq 1$.

For Zariski-generic $A^{(\alpha)}$, there exist $i_0,j_0$ such that $Q^{(\alpha\beta\gamma_0\delta_0)}_{i_0j_0k_0\ell_0} \neq 0$ for all distinct $(\alpha,\beta,\gamma_0,\delta_0)$. The rank-1 condition on $M^{\gamma_0\delta_0}_{k_0\ell_0}$ means there exist row-vectors $(r_\alpha^{(i)})$ and column-vectors $(c_\beta^{(j)})$ such that
\[
P^{(\alpha\beta\gamma_0\delta_0)}_{ij k_0\ell_0} = r_{(\alpha,i)} \cdot c_{(\beta,j)}.
\]

Dividing by the generically nonzero $Q^{(\alpha\beta\gamma_0\delta_0)}_{ijk_0\ell_0}$:
\[
\lambda_{\alpha\beta\gamma_0\delta_0} = \frac{r_{(\alpha,i)}c_{(\beta,j)}}{Q^{(\alpha\beta\gamma_0\delta_0)}_{ijk_0\ell_0}}.
\]

Since $\lambda_{\alpha\beta\gamma_0\delta_0}$ is independent of $i,j$, the ratio $r_{(\alpha,i)}c_{(\beta,j)}/Q^{(\alpha\beta\gamma_0\delta_0)}_{ijk_0\ell_0}$ must be independent of $i,j$. Using two values $i = i_0, i = i_1$ and $j = j_0$:
\[
\frac{r_{(\alpha,i_0)}}{Q^{(\alpha\beta\gamma_0\delta_0)}_{i_0j_0k_0\ell_0}} = \frac{r_{(\alpha,i_1)}}{Q^{(\alpha\beta\gamma_0\delta_0)}_{i_1j_0k_0\ell_0}}.
\]
By Zariski-genericity, $Q^{(\alpha\beta\gamma_0\delta_0)}_{ijk_0\ell_0} = \det[A^{(\alpha)}(i,:);A^{(\beta)}(j,:);A^{(\gamma_0)}(k_0,:);A^{(\delta_0)}(\ell_0,:)]$ depends on $i$ (generically). Define $u_\alpha = r_{(\alpha,i_0)}/Q^{(\alpha\beta_0\gamma_0\delta_0)}_{i_0j_0k_0\ell_0}$ for a fixed $\beta_0$. Then for all non-all-equal $(\alpha,\beta,\gamma_0,\delta_0)$:
\[
\lambda_{\alpha\beta\gamma_0\delta_0} = u_\alpha \cdot v_\beta,
\]
where $v_\beta = c_{(\beta,j_0)}/Q^{(\alpha_0\beta\gamma_0\delta_0)}_{i_0j_0k_0\ell_0}$ for a fixed $\alpha_0$.

By applying the vanishing of $\mathbf{F}$ for the other mode pairs ($(\alpha,\gamma)$-flattenings, $(\beta,\delta)$-flattenings, etc.), we similarly get $\lambda_{\alpha\beta\gamma\delta} = u_\alpha v_\beta w_\gamma x_\delta$ for all modes simultaneously. The consistency of these factorizations follows from the fact that $\mathbf{F} = 0$ imposes rank-1 conditions on ALL matrix flattenings of $\Lambda$, and a 4-way tensor is rank-1 iff all its matrix flattenings are rank-1 (since rank-1 of the $(\alpha,\beta)$-flattening gives $\lambda_{\alpha\beta\gamma\delta} = f_{\alpha\beta}\cdot g_{\gamma\delta}$, and rank-1 of the $(\alpha,\gamma)$-flattening gives $\lambda_{\alpha\beta\gamma\delta} = h_{\alpha\gamma}\cdot k_{\beta\delta}$; together these force the rank-1 product form $\lambda_{\alpha\beta\gamma\delta} = u_\alpha v_\beta w_\gamma x_\delta$).

\medskip
This completes the proof of Theorem~\ref{thm:P3}. The polynomial map $\mathbf{F}$ satisfies:
\begin{itemize}
\item \textbf{(P1):} $\mathbf{F}$ is defined purely in terms of $P^{(\alpha\beta\gamma\delta)} = \lambda_{\alpha\beta\gamma\delta}Q^{(\alpha\beta\gamma\delta)}$, with no explicit dependence on $A^{(\alpha)}$.
\item \textbf{(P2):} All coordinate functions have degree $2$, independent of $n$.
\item \textbf{(P3):} Proved in Theorem~\ref{thm:P3}.
\end{itemize}
\end{proof}

\section{Explicit Summary of $\mathbf{F}$}

The polynomial map $\mathbf{F}: \R^{81n^4} \to \R^N$ has the following explicit description:

For each of the $\binom{4}{2} = 6$ pairs of modes, for each choice of the remaining two mode-values $(\gamma,\delta) \in [n]^2$ (or equivalently the relevant fixed indices), for each choice of four distinct mode-indices $\alpha_1,\alpha_2,\beta_1,\beta_2 \in [n]$, and for each choice of local index tuples $(i_1,j_1,k_1,\ell_1)$ and $(i_2,j_2,k_2,\ell_2)$ in $[3]^4$, the map includes the coordinate function:
\[
P^{(\alpha_1\beta_1\gamma\delta)}_{i_1j_1k_1\ell_1} \cdot P^{(\alpha_2\beta_2\gamma\delta)}_{i_2j_2k_2\ell_2}
- P^{(\alpha_1\beta_2\gamma\delta)}_{i_1j_2k_1\ell_2} \cdot P^{(\alpha_2\beta_1\gamma\delta)}_{i_2j_1k_2\ell_1}.
\]
Each coordinate function is degree 2 in the input. The total number of coordinates $N = O(n^6 \cdot 3^8) = O(n^6)$ grows polynomially in $n$, but the \emph{degree} is fixed at 2.

\section{Verification Protocol}

\textbf{Check A (Completeness):} The $(\Leftarrow)$ direction: when $\lambda = u\otimes v\otimes w\otimes x$, the polynomial $F$ factors as $C \cdot (Q\text{-polynomial})$ with $C = u_{\alpha_1}u_{\alpha_2}v_{\beta_1}v_{\beta_2}w_\gamma^2 x_\delta^2 \neq 0$; the $Q$-polynomial vanishes by the Tucker structure. The $(\Rightarrow)$ direction: $F=0$ forces rank-1 of appropriate matrix flattenings of $\Lambda$; combined across all flattening pairs this gives $\lambda = u\otimes v\otimes w\otimes x$.

\textbf{Check B (Logic):} No circular reasoning. The Tucker structure (Lemma~\ref{lem:rescale}) is used in both directions. Genericity of $A^{(\alpha)}$ ensures $Q^{(\alpha\beta\gamma\delta)}_{ijk\ell} \neq 0$ for the reference indices needed.

\textbf{Check C (Definitions):} Property (P3) requires exactly the rank-1 factorization $\lambda_{\alpha\beta\gamma\delta} = u_\alpha v_\beta w_\gamma x_\delta$ for all non-all-equal indices, with $u,v,w,x \in (\R^*)^n$. The proof in the $(\Rightarrow)$ direction recovers this factorization.

\textbf{Check D (Generality):} The result holds for all $n \geq 5$ and Zariski-generic $A^{(\alpha)}$. Degrees are uniformly 2 independent of $n$.

\end{document}
